% This is a template for your written document.
%
% To compile using latexmk on the command line, run the following: 
% latexmk -pdf main.tex

\documentclass[12pt]{article}
\usepackage{setspace}
\singlespace
\usepackage[left=1in,right=1in,top=1in,bottom=1in]{geometry}
\usepackage{graphicx}

\title{\textbf{Project Topic}}
\author{Cameron Sentieri}

\begin{document}

\maketitle

\section{Introduction}
Within the field of Human-Computer Interaction (HCI), significant research has focused on designing systems, platforms, and content that capture and sustain user attention. Modern Smartphones employ persuasive design strategies, personalized algorithms, and continuous notifaction systems to maximize engagement. While these technologies have increased accessibility, communication, and productivity, they have also contributed to unprecedented levels of screen time and digital dependence. 

As engagement with digital devices become more constant, researchers have begun examining the broader implications of sustained technological interaction. Much of the existing literature explores impact on mental health productivity, and interpersonal relationships, However, comparatively less attention has been given to how these technologies influence spiritual life and religious practices, specifically in Christianity. From both psychological and theological perspecties, humans are often describes as relational beings whose well being depends on intentional connection with others (e.g., Genesis 2:18), and for many, with God. Practices such as prayer, solitude, worship, and reflection require sustained attention and cognitive presence, qualities that are challenged by distracting environments.

This paper investigates the question: How does technology affect our relationship with God? This will be done by a literature review that examines research in HCI, Customer Smartphone Distraction(CSD), Technology in spiritual disciplines, Prayer, and Scripture reading. By synthesizing these findings, this review identifies emerging themes, points of agreement, and areas where further research can be done.  
\section{Related Works}

\section{Methodology}
Dijkstra's algorithm was chosen for single-source shortest path computation on graphs with non-negative edge weights, following the implementation presented in CLRS~\cite{clrsAlgorithms}.

\section{Results}
An exemplary figure is shown in Figure~\ref{fig:dijkstra}. Any time a figure is used, it must have a capation and be directly referred to and discussed in the text.

%\begin{figure}[b] puts the image to the (closest) bottom of page
%\begin{figure}[t] puts the image to the (closest) top of page
%\begin{figure}[h] allows the image to float, putting it roughly wherever it was placed in the text.
\begin{figure}[b] 
\begin{center}
  \includegraphics[width=5cm]{imgs/dijkstra.png}
  \caption{Dijkstra's algorithm does not work on graphs with negative edge weights, as evidenced here.}
  \label{fig:dijkstra}
 \end{center}
\end{figure}



\section{Conclusion}

\bibliographystyle{acm}
\bibliography{bibliography.bib}

\end{document}
